\documentclass[10pt,a4paper]{amsart}
\usepackage{amsmath,amssymb,amsthm,graphicx}
\usepackage{verbatim}
\newenvironment{reference}{\comment}{\endcomment}
%\newenvironment{reference}{}{}
\newenvironment{slogan}{\comment}{\endcomment}
\newenvironment{history}{\comment}{\endcomment}

% For commutative diagrams we use Xy-pic
\usepackage[all]{xy}

% We use 2cell for 2-commutative diagrams.
\xyoption{2cell}
% Theorem environments.
%
\theoremstyle{plain}
\newtheorem{theorem}[subsection]{Theorem}
\newtheorem{proposition}[subsection]{Proposition}
\newtheorem{lemma}[subsection]{Lemma}

\theoremstyle{definition}
\newtheorem{definition}[subsection]{Definition}
\newtheorem{example}[subsection]{Example}
\newtheorem{exercise}[subsection]{Exercise}
\newtheorem{situation}[subsection]{Situation}

\theoremstyle{remark}
\newtheorem{remark}[subsection]{Remark}
\newtheorem{remarks}[subsection]{Remarks}

\numberwithin{equation}{subsection}

% Macros
%
\def\lim{\mathop{\mathrm{lim}}\nolimits}
\def\colim{\mathop{\mathrm{colim}}\nolimits}
\def\Spec{\mathop{\mathrm{Spec}}}
\def\Hom{\mathop{\mathrm{Hom}}\nolimits}
\def\Ext{\mathop{\mathrm{Ext}}\nolimits}
\def\SheafHom{\mathop{\mathcal{H}\!\mathit{om}}\nolimits}
\def\SheafExt{\mathop{\mathcal{E}\!\mathit{xt}}\nolimits}
\def\Sch{\mathit{Sch}}
\def\Mor{\mathop{\mathrm{Mor}}\nolimits}
\def\Ob{\mathop{\mathrm{Ob}}\nolimits}
\def\Sh{\mathop{\mathit{Sh}}\nolimits}
\def\NL{\mathop{N\!L}\nolimits}
\def\CH{\mathop{\mathrm{CH}}\nolimits}
\def\proetale{{pro\text{-}\acute{e}tale}}
\def\etale{{\acute{e}tale}}
\def\QCoh{\mathit{QCoh}}
\def\Ker{\mathop{\mathrm{Ker}}}
\def\Im{\mathop{\mathrm{Im}}}
\def\Coker{\mathop{\mathrm{Coker}}}
\def\Coim{\mathop{\mathrm{Coim}}}

% Boxtimes
%
\DeclareMathSymbol{\boxtimes}{\mathbin}{AMSa}{"02}

%
% Macros for moduli stacks/spaces
%
\def\QCohstack{\mathcal{QC}\!\mathit{oh}}
\def\Cohstack{\mathcal{C}\!\mathit{oh}}
\def\Spacesstack{\mathcal{S}\!\mathit{paces}}
\def\Quotfunctor{\mathrm{Quot}}
\def\Hilbfunctor{\mathrm{Hilb}}
\def\Curvesstack{\mathcal{C}\!\mathit{urves}}
\def\Polarizedstack{\mathcal{P}\!\mathit{olarized}}
\def\Complexesstack{\mathcal{C}\!\mathit{omplexes}}
% \Pic is the operator that assigns to X its picard group, usage \Pic(X)
% \Picardstack_{X/B} denotes the Picard stack of X over B
% \Picardfunctor_{X/B} denotes the Picard functor of X over B
\def\Pic{\mathop{\mathrm{Pic}}\nolimits}
\def\Picardstack{\mathcal{P}\!\mathit{ic}}
\def\Picardfunctor{\mathrm{Pic}}
\def\Deformationcategory{\mathcal{D}\!\mathit{ef}}
\usepackage{cleveref}
\crefname{originallemma}{original lemma}{original lemmas}
\Crefname{originallemma}{Original lemma}{Original lemmas}
\makeatletter
\newlabel{lemma-skyscraper-exact}{{05V7}{source}{}{sourcepointer-1}{}}
\newlabel{lemma-skyscraper-exact@cref}{{[originallemma][0][]05V7}{[1][0][]source}{}{}{}}
\makeatother
\crefname{originalreference}{original reference}{original references}
\Crefname{originalreference}{Original reference}{Original references}
\makeatletter
\newlabel{sites-section-example-sheaf-G-sets}{{00VZ}{source}{}{sourcepointer-2}{}}
\newlabel{sites-section-example-sheaf-G-sets@cref}{{[originalreference][0][]00VZ}{[1][0][]source}{}{}{}}
\makeatother
\crefname{originalreference}{original reference}{original references}
\Crefname{originalreference}{Original reference}{Original references}
\makeatletter
\newlabel{sites-example-point-G-sets}{{00YF}{source}{}{sourcepointer-3}{}}
\newlabel{sites-example-point-G-sets@cref}{{[originalreference][0][]00YF}{[1][0][]source}{}{}{}}
\makeatother
\title{Exactness of the skyscraper functor at a point of a ringed topos, and the non-linearity of its direct sum decomposition}
\author{The Stacks Project Authors}
\date{stacks-project:sites-modules.tex:local-snapshot-20261012; GNU FDL 1.2 or later}
\begin{document}
\maketitle
\noindent\textit{Editorial note.} All delivered mathematical text is the author's own, copied verbatim from the pinned Stacks Project chapter file sites-modules.tex; no statement or proof step was rewritten, reordered or removed.\par\medskip
\noindent\textit{Editorial note.} The counted span cites three external items (one lemma, one section, one example), all recorded in assembly.external\_references with canonical tags as navigation pointers.\par\medskip
\noindent\textit{Editorial note.} The author's italicised 'not' in part (3) is part of the source text and was copied verbatim.\par\medskip
\noindent\textit{Allowed standard prerequisites (editorial).}\begin{itemize}
\item Exactness and the direct sum decomposition for the abelian-sheaves version. Stacks Project, Modules on Sites, Lemma lemma-skyscraper-exact (tag 09B1). Supplies parts (1) and the surjectivity of (2), and the decomposition whose linearity fails in (3).
\item The topos of G-sets and its forgetful point. Stacks Project, Sites, Section sites-section-example-sheaf-G-sets (tag 00YR) and Example sites-example-point-G-sets (tag 00Z6). Provides the example on which the counterexample in (3) is built.
\end{itemize}\medskip
\ifdefined\hypertarget\hypertarget{sourcepointer-1}{}\fi
\noindent\textit{Source reference.} Original lemma 05V7: sources/stacks-chow/tags.txt, source lines 7541--7541. Cross-chapter standard prerequisite cited by the author inside the counted span; it is used as a pointer and its own proof is not part of this item.\par\medskip
\ifdefined\hypertarget\hypertarget{sourcepointer-2}{}\fi
\noindent\textit{Source reference.} Original reference 00VZ: sources/stacks-chow/tags.txt, source lines 1165--1165. Cross-chapter standard prerequisite cited by the author inside the counted span; it is used as a pointer and its own proof is not part of this item.\par\medskip
\ifdefined\hypertarget\hypertarget{sourcepointer-3}{}\fi
\noindent\textit{Source reference.} Original reference 00YF: sources/stacks-chow/tags.txt, source lines 1253--1253. Cross-chapter standard prerequisite cited by the author inside the counted span; it is used as a pointer and its own proof is not part of this item.\par\medskip
\begin{lemma}
\label{lemma-skyscraper-modules-exact}
Let $(\Sh(\mathcal{C}), \mathcal{O})$ be a ringed topos.
Let $p$ be a point of the topos $\Sh(\mathcal{C})$.
\begin{enumerate}
\item The functor
$p_* : \textit{Mod}(\mathcal{O}_p) \to \textit{Mod}(\mathcal{O})$,
$M \mapsto p_*M$ is exact.
\item The canonical surjection $p^{-1}p_*M \to M$ is $\mathcal{O}_p$-linear.
\item The functorial direct sum decomposition
$p^{-1}p_*M = M \oplus I(M)$ of Lemma \ref{lemma-skyscraper-exact}
is {\bf not} $\mathcal{O}_p$-linear in general.
\end{enumerate}
\end{lemma}

\begin{proof}
Part (1) and surjectivity in (2)
follow immediately from the corresponding result for abelian
sheaves in
Lemma \ref{lemma-skyscraper-exact}.
Since $p^{-1}\mathcal{O} = \mathcal{O}_p$
we have $p^{-1} = p^*$ and hence $p^{-1}p_*M \to M$
is the same as the counit $p^*p_*M \to M$
of the adjunction for modules, whence linear.

\medskip\noindent
Proof of (3). Suppose that $G$ is a group. Consider the topos
$G\textit{-Sets} = \Sh(\mathcal{T}_G)$
and the point $p : \textit{Sets} \to G\textit{-Sets}$.
See Sites, Section \ref{sites-section-example-sheaf-G-sets} and
Example \ref{sites-example-point-G-sets}.
Here $p^{-1}$ is the functor forgetting
about the $G$-action. And $p_*$ is the right adjoint
of the forgetful functor, sending $M$ to $\text{Map}(G, M)$.
The maps in the direct sum decomposition are the maps
$$
M \to \text{Map}(G, M) \to M
$$
where the first sends $m \in M$ to the constant map with value $m$ and
where the second map is evaluation at the identity element $1$ of $G$.
Next, suppose that $R$ is a ring endowed with an action of $G$.
This determines a sheaf of rings $\mathcal{O}$ on $\mathcal{T}_G$.
The category of $\mathcal{O}$-modules is the category of $R$-modules $M$
endowed with an action of $G$ compatible with the action on $R$.
The $R$-module structure on $\text{Map}(G, M)$ is given by
$$
( r f ) (\sigma) = \sigma(r) f(\sigma)
$$
for $r \in R$ and $f \in \text{Map}(G, M)$. This is true because it is the
unique $G$-invariant $R$-module structure compatible with evaluation at $1$.
The reader observes that in general the image of $M \to \text{Map}(G, M)$
is not an $R$-submodule (for example take $M = R$ and assume the
$G$-action is nontrivial), which concludes the proof.
\end{proof}

\end{document}
